% Parte I: copia de docs/INFORME_FINAL.md. Mismo texto y mismas cifras; solo se agregan
% referencias cruzadas a las partes II y III.

\section{Resumen}\label{sec:resumen}

\begin{itemize}
  \item \textbf{Qué es.} Alicanto es un proxy que se pone entre un agente de código y sus herramientas.
    En cada llamada mide cuántos bits de cada clave ya salieron, aunque salgan de a cuatro caracteres o
    codificados, y le pide al gestor de secretos que cambie la clave antes de que el atacante la
    complete. No modifica el agente y ningún LLM decide.
  \item \textbf{Qué quedó hecho.} El producto corre completo en local: proxy, motor de medición, gestor
    de secretos con rotación, API protegida y servidor del atacante. Tiene un \esp{hook} para Claude
    Code, una demo desplegada en \href{https://alicanto.vercel.app}{alicanto.vercel.app}, tres demos
    alternativas y un pitch de 3 minutos con una parte que corre en vivo. Hay más de 800 tests
    automáticos.
  \item \textbf{Cómo se midió.} Con tres roles separados. El equipo rojo escribió los ataques sin abrir
    el detector y selló el set de prueba con un hash. El equipo azul congeló el detector antes de que
    el set se abriera. Las reglas del equipo morado prohibían ajustar mirando la prueba, y la
    evaluación se corrió una sola vez.
  \item \textbf{El resultado.} De 180 ataques grabados, de seis familias que el detector no había
    visto, Alicanto frenó 120 tal como los ve el proxy en vivo: cuatro de las seis familias. No bloqueó
    ninguna de \num{5879} llamadas normales.
  \item \textbf{Lo que cuesta y lo que no cubre.} Rotó una clave de más en 34 de 90 sesiones normales
    difíciles. Dos familias se escapan: la clave enviada a su propio destino permitido y la fuga por
    \esp{commits} de git. Las dos están publicadas, con el comando para repetir la medición. Son
    ataques grabados y tráfico sintético, no un atacante que se adapta ni tráfico de producción.
  \item \textbf{Lo que se pide.} Dos equipos que trabajen con agentes de código, para un piloto.
\end{itemize}

\section{El problema y la idea}\label{sec:problema}

Un agente de código puede sacar una clave del repositorio en pedazos tan chicos que ninguna llamada
parece un robo. Alicanto mide cuánta clave ya salió y la cambia antes de que el atacante la complete.

\begin{itemize}
  \item \textbf{Por qué pasa.} Los agentes leen texto que escribe cualquiera: issues, páginas,
    documentación. Una instrucción plantada ahí puede pedirles que envíen una clave de a cuatro
    caracteres, entre decenas de llamadas normales.
  \item \textbf{Por qué no basta lo habitual.} Un guardia que mira cada llamada por separado, o cada
    agente por separado, no ve la suma. Si el atacante reparte la clave entre varios agentes, cada uno
    queda bajo su propio umbral.
  \item \textbf{La idea.} No adivinamos la intención. Contamos en bits cuánto de cada clave ya salió,
    sumando entre llamadas, codificaciones y agentes, y la rotamos en el gestor de secretos antes de
    que se complete.
  \item \textbf{Por qué rotar y no bloquear.} Bloquear interrumpe trabajo legítimo cuando el guardia se
    equivoca. Rotar no detiene a nadie: deja sin valor lo que ya salió.
  \item \textbf{Un precedente público.} En 2025, CamoLeak~\cite{camoleak} sacó secretos de
    repositorios privados a través de GitHub Copilot Chat, carácter por carácter. Es un caso parecido,
    no el mismo escenario de nuestra demo.
\end{itemize}

\section{Lo que construimos}\label{sec:construimos}

Alicanto se instala donde ya pasan las llamadas del agente. Es un servidor MCP que le ofrece al agente
sus herramientas (leer archivos, listar carpetas, peticiones HTTP, git, issues y comentarios de GitHub)
y decide sobre cada llamada antes de dejarla salir.

\begin{table}[ht]
\caption{Las piezas de Alicanto. La parte II explica cada una.}\label{tab:piezas}
\small
\begin{tabularx}{\textwidth}{@{}P{3.1cm}LP{3.6cm}@{}}
\toprule
Pieza & Qué hace & Dónde está \\
\midrule
Proxy MCP & Ofrece seis herramientas al agente y pasa cada llamada por el detector. Si el gestor de
  secretos no responde, bloquea las salidas & \cod{proxy/} \\
Contador de bits & Lleva la cuenta de cuántos bits de cada clave ya salieron. Lee fragmentos desde 4
  caracteres, en claro, base64, hexadecimal, URL, ROT13, invertidos y separados & \cod{engine/bits/} \\
Evidencia de comportamiento & Un modelo simple (histogramas y una suma acumulada) sobre señales de
  cada llamada: si la sesión leyó un secreto, hacia dónde sale la llamada & \cod{engine/evidence/} \\
Decisión & Permite, bloquea o rota. Tiene un tope de tres rotaciones cada 200 llamadas; después
  bloquea las salidas de esa clave y avisa a una persona & \cod{engine/detector/} \\
Cuenta compartida & Una sola cuenta por clave para todos los agentes de una misma máquina: una fuga
  repartida entre varios agentes se suma & \cod{engine/detector/ledger.py} \\
Servicios & Gestor de secretos con rotación, API protegida que responde 200 o 401, servidor del
  atacante que cuenta lo que recibió, y GitHub y web simulados & \cod{services/} \\
\esp{Hook} de Claude Code & Pasa \cod{Bash}, \cod{WebFetch} y \cod{Read} por el contador, sin cambiar
  las herramientas del agente & \cod{integrations/claude_code/} \\
\bottomrule
\end{tabularx}
\end{table}

Dos decisiones de diseño explican el resto.

\textbf{Medir en vez de adivinar intenciones.} Los guardias habituales miran cada llamada por separado
y deciden si parece un robo. Cuatro caracteres por llamada nunca lo parecen. Alicanto no juzga la
llamada: suma lo que ya salió de cada clave.

\textbf{Rotar en vez de bloquear.} Cuando la cuenta llega al umbral, la clave se cambia en el gestor de
secretos. Lo que el atacante juntó es de una clave que ya no sirve, y el agente sigue trabajando. Un
falso positivo cuesta una rotación, no una tarea interrumpida.

La clave de las pruebas tiene 160 bits y el ataque triunfa si saca más de 96. Hay dos configuraciones,
y conviene no confundirlas. La demo rota cuando salió el \qty{30}{\percent} de la clave (48 bits) y
corre solo con el contador de bits, para que se vea el goteo. La evaluación usa la configuración
congelada, que rota al \qty{10}{\percent} y suma la evidencia de comportamiento.

\section{Cómo lo evaluamos: equipo rojo, azul y morado}\label{sec:evaluamos}

El riesgo de evaluar un detector que uno mismo construye es ajustarlo hasta que pase su propia prueba.
Para evitarlo separamos tres roles. No son tres grupos de personas: son tres papeles, con reglas
escritas antes de empezar (\cod{docs/PLAN.md}).

\begin{table}[ht]
\caption{Los tres roles de la evaluación.}\label{tab:roles}
\small
\begin{tabularx}{\textwidth}{@{}P{2.4cm}LL@{}}
\toprule
Rol & Qué hizo & Qué lo ata \\
\midrule
\textcolor{rojoeq}{\rule{0.55em}{0.55em}}\ Equipo rojo & Escribió los ataques y el tráfico normal
  difícil, en una sesión de agente aparte & No abre el código del detector hasta el congelamiento \\
\textcolor{azuleq}{\rule{0.55em}{0.55em}}\ Equipo azul & Construyó el detector y eligió sus parámetros
  con los sets de entrenamiento y validación & No agrega reglas para un ataque concreto. Cada cambio
  queda en una bitácora, con hora y motivo \\
\textcolor{moradoeq}{\rule{0.55em}{0.55em}}\ Equipo morado & Puso las reglas y las hizo cumplir & El
  set de prueba se sella con hash, el detector se congela con hash y la evaluación se corre una sola
  vez \\
\bottomrule
\end{tabularx}
\end{table}

\subsection{Los tres sets}

\begin{table}[ht]
\caption{Los tres sets de datos.}\label{tab:sets}
\small
\begin{tabularx}{\textwidth}{@{}P{3.2cm}R{1.5cm}R{3.2cm}L@{}}
\toprule
Set & Ataques & Sesiones normales difíciles & Familias \\
\midrule
Entrenamiento & 80 & 20 & A1, A2, A3, A4, A5$'$; B1, B2 \\
Validación & 80 & 20 & Las mismas \\
Prueba (sellado) & 180 & 90 & A5, A7, A9, A10, A12, A13; B3, B4, B5 \\
\bottomrule
\end{tabularx}
\end{table}

Las familias del set de prueba no aparecen en los otros dos: usan canales y codificaciones nuevos.
Cada familia tiene 30 casos generados a partir de una misma plantilla; en A7, cada caso es un grupo de
tres sesiones.

\subsection{La secuencia del sábado 21 de noviembre}

\begin{enumerate}
  \item \textbf{07:10.} El rojo sella el set de prueba: publica el sha256 de todos sus archivos,
    \cod{18dda3af}…, en \cod{redteam/SELLO.md}. El set queda en una rama aparte.
  \item \textbf{Antes de abrir.} El azul congela el detector. Parámetros y modelo quedan fijados con el
    hash \cod{cf56ce20}… (commit \cod{3f26307}, etiqueta \cod{congelamiento}).
  \item \textbf{09:40.} Se abre el set. El sello se recalcula y coincide. La evaluación se corre una
    sola vez (commit \cod{cd14e9e}); el comando se niega a correr de nuevo.
  \item \textbf{Después.} Los análisis posteriores (el atacante que se adapta, los fragmentos cortos,
    las marcas del proxy) se hicieron con el detector congelado y se reportan aparte. Nada se ajustó
    después de ver los resultados.
\end{enumerate}

\subsection{Qué puede comprobar alguien de afuera}

\begin{itemize}
  \item \textbf{El sello.} El comando de \cod{redteam/SELLO.md} vuelve a calcular el hash sobre los
    archivos de hoy, y coincide. La demo del pitch lo recalcula en vivo.
  \item \textbf{El orden.} El sello se publicó antes del congelamiento y de la corrida; los commits
    están en el historial.
  \item \textbf{Las cifras.} \cod{python -m eval.final.proxy_flags --out /tmp/pf} repite la medición en
    unos tres minutos. La repetimos el 21 de noviembre a las 19:20 y dio las mismas cifras.
\end{itemize}

Lo que no puede comprobar es que el rojo no abrió el detector. Es una regla de trabajo del equipo, y lo
decimos así.

\section{Resultados}\label{sec:resultados}

\subsection{La cifra principal: 120 de 180}

La evaluación se corrió primero con las marcas que traían las grabaciones. Cada llamada lleva dos: si
saca datos y si su contenido es no confiable. Las había escrito el generador de ataques, y con ellas
Alicanto previno 180 de 180. Horas después el equipo encontró que el proxy en vivo calcula esas marcas
de otra forma en 810 de \num{15281} llamadas. Con las marcas del proxy, que es como funcionaría
instalado, la cifra es 120 de 180. Esa es la que decimos.

\begin{table}[ht]
\caption{Ataques frenados por familia, con las marcas del proxy (\cod{results/proxy_flags/systems.md}).
El detalle de por qué se frena o se escapa cada una está en la \cref{sec:por-familia}.}\label{tab:familias-i}
\small
\begin{tabularx}{\textwidth}{@{}P{1.4cm}LR{2.6cm}@{}}
\toprule
Familia & Cómo saca la clave & Frenados, de 30 \\
\midrule
A5 & Por un canal de confianza del propio entorno, en base32 & 30 \\
A7 & Repartida entre tres agentes; ninguno saca más de un tercio & 30 \\
A10 & Ofuscada con XOR y una clave por sesión & 30 \\
A12 & Goteo lento: decenas de llamadas normales y después una ráfaga & 30 \\
A9 & Hacia su propio destino permitido, en una cabecera & 0 \\
A13 & En mensajes de \esp{commits} de git, y después un \esp{push} & 0 \\
\bottomrule
\end{tabularx}
\end{table}

Cada familia son 30 réplicas de una misma plantilla. El resultado se lee como «cuatro de seis
familias», no como 180 casos independientes.

\textbf{Por qué se escapan dos.} En A9 la clave viaja a un destino que tiene permitido: para el proxy
ese envío es legítimo y no lo cuenta como fuga. En A13 el \cod{git commit} no es una salida; la salida
es el \cod{push}. El proxy en vivo inspecciona el \esp{push}, rota y lo bloquea en las 30 sesiones,
pero el set marca la fuga en cada \esp{commit}, así que la métrica da 0 de 30. La métrica no se tocó.

\subsection{El tráfico normal}

Ninguna de \num{5879} llamadas normales fue bloqueada. El costo está en las rotaciones: en 34 de las
90 sesiones normales difíciles Alicanto rotó una clave sin necesidad (39 rotaciones en total). Casi
todas son de una sola familia.

\begin{table}[ht]
\caption{Sesiones normales difíciles en que Alicanto rotó una clave sin necesidad
(\cod{results/test/summary_ours.json}).}\label{tab:normales-i}
\small
\begin{tabularx}{\textwidth}{@{}LR{4.2cm}@{}}
\toprule
Sesiones normales difíciles & Con una rotación de más, de 30 \\
\midrule
B3: identificadores que parecen claves (SHA, UUID, tokens) & 3 \\
B4: despliegue y rotación legítimos & 1 \\
B5: contenido no confiable lleno de claves de ejemplo & 30 \\
\bottomrule
\end{tabularx}
\end{table}

B5 hace, de forma legítima, lo mismo que un ataque: lee un secreto y después habla con un destino no
confiable. En las 30 sesiones normales corrientes no rotó nunca.

\subsection{La misma medición con otros guardias}

Con las marcas del proxy, sobre los mismos 180 ataques:

\begin{table}[ht]
\caption{La misma medición con otros sistemas (\cod{results/proxy_flags/systems.md}). La
\cref{tab:sistemas} trae los intervalos y las demás columnas.}\label{tab:guardias-i}
\small
\begin{tabularx}{\textwidth}{@{}LR{2.6cm}R{3.4cm}@{}}
\toprule
Sistema & Ataques frenados & Llamadas normales bloqueadas \\
\midrule
Alicanto & 120 & 0 \\
Regla simple: «leyó un secreto y sale a un destino no confiable» & 90 & 0 \\
Presupuesto de bits por sesión (réplica nuestra, simplificada) & 40 & \qty{4.8}{\percent} \\
Lista de destinos permitidos & 30 & 0 \\
Lista de destinos más escáner de secretos & 30 & \qty{2.2}{\percent} \\
Alicanto sin evidencia de comportamiento & 10 & 0 \\
Alicanto sin rotación & 0 & 0 \\
Escáner de secretos (regex y entropía) & 0 & \qty{2.2}{\percent} \\
Reensamblado por sesión (réplica nuestra, simplificada) & 0 & 0 \\
\bottomrule
\end{tabularx}
\end{table}

Las dos réplicas son versiones simplificadas que hicimos a partir de documentación pública. Sus cifras
no son las de ningún producto.

\subsection{Mediciones posteriores, con el detector congelado}

\begin{itemize}
  \item \textbf{Atacante que se adapta.} Es una simulación, con las marcas de las grabaciones: el
    atacante relee la clave nueva después de cada rotación. Ninguno de los 180 pasa de 96 bits de una
    versión de la clave. En todos los casos Alicanto llega a su tope de tres rotaciones y bloquea las
    salidas de esa clave.
  \item \textbf{Fragmentos cortos.} Con pedazos de 4 y 7 caracteres, ninguna de 36 combinaciones pasa
    de 96 bits. Los pedazos de 1 a 3 caracteres el contador no los cuenta: la configuración evaluada
    los frena por la evidencia de comportamiento (máximo 12 bits), pero la de la demo y del
    \esp{hook}, sin ese modelo, deja salir la clave entera en 10 de 18 combinaciones.
  \item \textbf{Modelos reales.} Un agente con Claude Haiku~5.5, Sonnet~5.5 y Opus~5.5 trabajó a
    través del proxy. En 21 corridas con la nota plantada ninguno la siguió, así que no hubo fuga que
    frenar. Lo que estas corridas muestran es que un agente con un modelo real trabaja a través del
    proxy; en una de ellas Alicanto rotó de más.
\end{itemize}

\section{Límites}\label{sec:limites-i}

Van aquí, juntos, para que nadie tenga que buscarlos. La \cref{sec:limites} los amplía con el detalle
técnico.

\begin{itemize}
  \item \textbf{Son ataques grabados.} Las sesiones del set son guiones: después de una rotación
    siguen enviando la clave vieja. El atacante que se adapta se midió aparte, y en simulación.
  \item \textbf{Todo el tráfico es sintético.} Ataques y sesiones normales los generó el equipo. No
    hay tráfico de producción.
  \item \textbf{Dos familias no se cubren:} la clave enviada a su propio destino permitido (A9) y la
    fuga por \esp{commits} de git (A13).
  \item \textbf{Lo que frena los ataques nuevos es la evidencia de comportamiento, no el contador.}
    Solo con el contador de bits se frenan 10 de 180. El contador es exacto para las codificaciones
    que conoce y no detecta las que no conoce.
  \item \textbf{El modelo de evidencia equivale, en estos datos, a una regla simple:} «la sesión leyó
    un secreto y esta llamada sale a un destino no confiable». Con las marcas de las grabaciones esa
    regla empata con Alicanto. Con las del proxy, Alicanto frena 120 y la regla 90.
  \item \textbf{Rotar de más tiene un costo que no cumple nuestra propia meta:} 5,46 rotaciones por
    cada \num{1000} llamadas normales, cuando la meta era 5.
  \item \textbf{El contador no cuenta pedazos de menos de 4 caracteres.}
  \item \textbf{La cuenta compartida cubre varios agentes en una misma máquina,} no varias máquinas.
  \item \textbf{Rotar no impide el intento.} Hace que lo robado no sirva.
  \item \textbf{El agente de las demos es un guion.} No mostramos a Alicanto frenando a un modelo real,
    porque ningún modelo siguió la nota plantada.
  \item \textbf{El \esp{hook} de Claude Code ve lo que el modelo escribe en los comandos,} no lo que
    un comando hace por dentro. Se probó en una sesión real; el contador con un modelo real no quedó
    probado.
  \item \textbf{No hay evidencia de demanda todavía.} Ningún equipo externo lo ha usado. Por eso el
    pedido es un piloto.
\end{itemize}

\section{La demo y el pitch}\label{sec:demo-i}

\subsection{La demo oficial}

Está en \cod{demo/viz/} y se corre en local. Fue la que publicó el sitio hasta el domingo 22 a las
05:27. Desde entonces \href{https://alicanto.vercel.app}{alicanto.vercel.app} publica la alternativa
03, con sus terminales, el tablero de pruebas y este informe.
Muestra el mismo ataque dos veces. Sin protección, el atacante arma la clave y recibe 200. Con
Alicanto, la clave rota en la llamada 24, con 48 bits filtrados, y recibe 401. Una segunda versión
reparte la fuga entre cinco agentes: con una cuenta por agente nadie la ve; con la cuenta por clave,
rota.

\subsection{Las demos alternativas}

Son tres exploraciones en \cod{demo/alternativas/}, hechas sobre corridas reales del motor.

\begin{table}[ht]
\caption{Las demos alternativas.}\label{tab:alternativas}
\small
\begin{tabularx}{\textwidth}{@{}P{4.6cm}L@{}}
\toprule
Alternativa & Qué muestra \\
\midrule
01, «El botín» & El ataque visto desde el servidor del atacante: la clave se arma de a cuatro
  caracteres y, con Alicanto, lo que había juntado se vuelve ceniza \\
02, «El juez es el atacante» & Una página para el teléfono de alguien del jurado, que recibe la clave
  pedazo a pedazo y la prueba \\
03, «20 agentes, una clave» & Una clínica ficticia, sus 20 asistentes de código y una nota plantada.
  Es la que lleva el pitch de 3 minutos y, desde el domingo, la del sitio público \\
\bottomrule
\end{tabularx}
\end{table}

\subsection{El pitch de 3 minutos}

Con \cod{?pitch=1}, la alternativa 03 lleva los 3 minutos completos en una sola página, con reloj.

\begin{table}[ht]
\caption{El pitch de 3 minutos, por tramos.}\label{tab:pitch}
\small
\begin{tabularx}{\textwidth}{@{}P{2.2cm}LP{3.9cm}@{}}
\toprule
Reloj & Qué pasa & De dónde sale \\
\midrule
0:00 a 0:40 & Apertura: qué es Alicanto y por qué importa & Texto \\
0:40 a 1:23 & El ataque. Una nota plantada en un issue; diez asistentes mandan cuatro caracteres cada
  uno; ningún monitor por agente ve más de 16 de 160 bits; el atacante arma la clave y abre una ficha
  & Corrida grabada del motor \\
1:23 a 1:51 & El público hace de atacante: elige entre cuántos asistentes reparte la clave, cómo
  esconde cada pedazo y un número cualquiera & El motor corre en ese momento, en unos 200~ms \\
1:51 a 2:21 & Con Alicanto, frente al ataque que eligieron: la cuenta llega a 48 bits, la clave rota,
  el atacante recibe 401 y ningún asistente queda bloqueado & Esa misma corrida \\
2:21 a 2:53 & Cómo lo probamos: los tres equipos, el set sellado, 120 de 180, sus límites y el pedido
  & Cifras de \cod{results/}; el sello se recalcula en vivo \\
\bottomrule
\end{tabularx}
\end{table}

Los tiempos son de un ensayo a velocidad real, con el guion dicho a 150 palabras por minuto. Terminó
en 2:53.

\textbf{La última escena} responde a «¿y esto funciona de verdad?». Dibuja el set de prueba entero
como una grilla: una fila por familia y una casilla por ataque o por sesión normal. Primero aparece
sellado, con su hash. Después se abre y las casillas se llenan: claras las de los ataques frenados, en
oro las de las dos familias que se escapan y en gris las sesiones normales en que rotó de más. Un paso
más, para las preguntas, muestra la misma medición con otros guardias. La página no trae ninguna cifra
escrita: las lee de un archivo que un script copia de \cod{results/}, y un test falla si dejan de
coincidir. La \cref{fig:rejilla} de este documento es esa misma grilla.

\textbf{Qué es real y qué no.} El motor, sus decisiones, la rotación y el resultado del cobro son
reales. La pantalla de la clínica, la tarea que se escribe y la ficha son maquetas, y la página lo
dice. El agente es un guion.

El guion (\cod{GUION_3_MIN.md}) es un borrador. Lo tiene que reescribir, con su voz, quien presente.

\section{Cómo trabajamos}\label{sec:trabajamos}

Cinco personas, cada una con su agente de código, durante 36 horas. El trabajo se repartió por áreas
con dueño, y los agentes se coordinaron entre sí por escrito.

\begin{table}[ht]
\caption{Las áreas del equipo, según \cod{AGENTS.md}.}\label{tab:areas}
\small
\begin{tabularx}{\textwidth}{@{}P{2.3cm}LP{2.9cm}P{3.9cm}@{}}
\toprule
Área & Qué cubre & Dueño & Carpetas \\
\midrule
Líder y orquestador & Plan, reglas, contratos compartidos y reparto de tareas & \cod{jplira2} &
  \cod{docs/PLAN.md}, \cod{AGENTS.md}, \cod{contracts/} \\
Proxy & Proxy MCP, herramientas y grafos & \cod{jcbasulto} & \cod{proxy/}, \cod{engine/graph/} \\
Análisis & Motor de bits, evidencia, búsqueda de parámetros y evaluación & \cod{matix3000} &
  \cod{engine/}, \cod{eval/} \\
Servicios & Gestor de secretos, API protegida, servidor del atacante y líneas de base &
  \cod{terdmannsdorffer} & \cod{services/}, \cod{eval/baselines/} \\
Rojo & Ataques, set sellado, GitHub y web simulados & \cod{jplira2} & \cod{redteam/} \\
Demo & Demo, agente mínimo, visualización y pitch & \cod{VictoriaGuerriero} & \cod{demo/},
  \cod{pitch/} \\
\bottomrule
\end{tabularx}
\end{table}

Las áreas son las de \cod{AGENTS.md}. Durante el evento el líder movió trabajo entre personas más de
una vez; el detalle está en \cod{docs/INFORME_TURNOS.md}.

\subsection{Cómo se coordinaron los agentes}

\begin{itemize}
  \item \textbf{Un tablero para agentes.} Un issue fijado de GitHub hace de tablero. Cada agente
    publica en una línea qué cambió, qué contrato tocó y qué necesita de quién, y lee lo de los demás
    antes de cada commit (\cod{tools/board}).
  \item \textbf{Un orquestador.} Un agente del líder reparte el plan como tareas. Cada tarea es un
    issue con objetivo, archivos, contrato y un comando de aceptación. Nadie reporta «listo» sin haber
    visto pasar ese comando.
  \item \textbf{Un PR por tarea.} Ramas cortas, commits solo con archivos propios, y aviso en el
    tablero cuando un merge cambia algo de lo que dependen otros.
  \item \textbf{Contratos compartidos.} Los tipos comunes (\cod{ToolCall}, \cod{Decision}, el registro
    de eventos) viven en \cod{contracts/} y solo los cambia el líder.
  \item \textbf{Reglas contra el sobreajuste,} descritas más arriba, y una bitácora de cada cambio al
    detector.
\end{itemize}

En números, en la rama principal al sábado 21 a las 17:40: 167 commits, 109 PR mergeados, 83 issues
cerrados y más de 800 tests automáticos.

\subsection{La línea de tiempo}

\begin{table}[ht]
\caption{La línea de tiempo del evento.}\label{tab:tiempo}
\small
\begin{tabularx}{\textwidth}{@{}P{3.6cm}L@{}}
\toprule
Cuándo & Qué pasó \\
\midrule
Viernes 20, 23:00 & Parte el hackathon \\
Sábado, 02:00 a 04:00 & Producto base en la rama principal y primera demo de punta a punta: 200 sin
  protección, 401 con Alicanto \\
04:00 a 06:00 & Escenario, ataques de entrenamiento y validación, primer nombre (Cuentagotas) \\
06:00 a 08:00 & Cuenta compartida entre agentes. El rojo sella el set de prueba (07:10) \\
08:00 a 09:30 & Auditoría del código contra lo que afirmábamos \\
09:40 & Congelamiento del detector y evaluación final, una sola vez \\
09:55 a 11:00 & Atacante que se adapta, medido aparte. Demo desplegada \\
11:30 & El proyecto pasa a llamarse Alicanto. Corridas con modelos reales \\
13:00 a 14:00 & Hallazgo de las marcas: la cifra pasa de 180 a 120 \\
14:00 a 16:25 & \esp{Hook} de Claude Code, fragmentos cortos, demos alternativas 01 y 02 \\
Tarde y noche & Alternativa 03, pitch de 3 minutos con la parte en vivo y la escena «Cómo lo
  probamos» \\
\bottomrule
\end{tabularx}
\end{table}

\subsection{Lo que corregimos en el camino}

Varias de las cosas más útiles del proyecto salieron de buscarle fallas.

\begin{itemize}
  \item \textbf{El motor no veía fragmentos de cuatro caracteres en base64.} Se encontró y se
    corrigió antes de congelar.
  \item \textbf{Dos afirmaciones nuestras eran falsas.} Decíamos que nadie juntaba fragmentos entre
    llamadas; una investigación del estado del arte mostró que ya hay herramientas que lo hacen. Se
    corrigieron la descripción y el pitch.
  \item \textbf{La evaluación contaba mal las rotaciones} y la herramienta \cod{git} del proxy tenía
    un hueco de seguridad. Los encontró una auditoría del código; los dos se corrigieron antes del
    congelamiento.
  \item \textbf{El 180 de 180 dependía de marcas que escribió el generador de ataques.} Lo encontró
    una revisión crítica contra la rúbrica. La decisión del líder fue recalcular y publicar la cifra
    con las marcas del proxy (120), en vez de solo agregar una salvedad. Hoy hay tests que vigilan que
    ningún texto público diga «180 de 180» sin su explicación.
  \item \textbf{La frase de los equipos separados} pasó a decir solo lo que se puede sostener: una
    sesión aparte con una regla de trabajo, no algo comprobable desde fuera.
\end{itemize}

\section{Qué queda pendiente}\label{sec:pendiente}

\subsection{Antes de la entrega del domingo a las 09:30}

\begin{itemize}
  \item \textbf{Elegir la demo del pitch.} La oficial o la alternativa 03 con su guion de 3 minutos.
    Lo decide el líder.
  \item \textbf{Reescribir y ensayar el guion.} Es un borrador. Quien presente lo dice con su voz y lo
    pasa con reloj, en el proyector de verdad: tiene que terminar antes de 2:55.
  \item \textbf{Revisar \cod{README.md} y \cod{project-description.md}.} Los preparó un agente; una
    persona del equipo los reescribe y los revisa antes de entregar. El PR \#178, con la descripción
    reescrita, espera al líder.
  \item \textbf{Si se elige la alternativa 03,} llevarla al sitio publicado y redesplegar. Hoy se abre
    solo en local.
\end{itemize}

\subsection{Lo que no alcanzamos a medir}

\begin{itemize}
  \item \textbf{Un atacante real que se adapta.} Lo medido es una simulación con un guion que relee y
    reenvía.
  \item \textbf{El costo del modo escalado sobre tráfico normal.} Cuando Alicanto llega a su tope de
    rotaciones y bloquea las salidas de una clave, no sabemos cuánto molesta a un equipo que trabaja.
  \item \textbf{Tráfico de producción.} Para eso es el piloto.
  \item \textbf{Varias máquinas.} La cuenta compartida funciona entre agentes de una misma máquina.
\end{itemize}

\subsection{Qué haríamos después, si hay piloto}

\begin{itemize}
  \item Cubrir las dos familias que hoy se escapan: la clave que viaja a su propio destino permitido
    y la fuga por \esp{commits} de git.
  \item Bajar las rotaciones de más en sesiones que leen contenido no confiable.
  \item Probarlo con dos equipos que usen agentes de código y tengan claves en sus repositorios, y
    medir con su tráfico.
\end{itemize}

\section{Dónde está cada cosa}\label{sec:donde}

Todo está en el repositorio del equipo, \cod{platanus-hack/platanus-hack-26-cl-team-47}.

\begin{table}[ht]
\caption{Dónde está cada cosa en el repositorio. El \apref{ap:mapa} trae el mapa completo.}\label{tab:donde}
\small
\begin{tabularx}{\textwidth}{@{}LP{7.4cm}@{}}
\toprule
Qué & Dónde \\
\midrule
Qué es el proyecto y cómo correrlo & \cod{README.md}, \cod{docs/QUICKSTART.md} \\
El plan y las reglas del equipo & \cod{docs/PLAN.md}, \cod{AGENTS.md} \\
La historia, turno por turno & \cod{docs/INFORME_TURNOS.md}, \cod{docs/ESTADO.md} \\
Cómo leer los resultados y qué frases respaldan & \cod{results/test/LECTURA.md} \\
La cifra con las marcas del proxy & \cod{results/proxy_flags/systems.md} \\
El sello del set de prueba & \cod{redteam/SELLO.md} \\
Las familias de ataque y sus fuentes & \cod{redteam/README.md}, \cod{redteam/test/FAMILIAS.md} \\
Cada cambio al detector, con hora y motivo & \cod{docs/BITACORA_DETECTOR.md} \\
El atacante que se adapta y los fragmentos cortos & \cod{results/adaptive/}, \cod{results/short_fragments/} \\
Las corridas con modelos reales & \cod{docs/LLM_EN_VIVO.md} \\
El \esp{hook} de Claude Code & \cod{integrations/claude_code/README.md} \\
La demo oficial & \cod{demo/viz/} \\
El sitio público y su despliegue & \cod{demo/alternativas/03-enjambre-juanpedro/publicar.sh}, \cod{vercel.json}, \cod{demo/DEPLOY.md} \\
Las demos alternativas & \cod{demo/alternativas/} \\
El pitch de 3 minutos con la parte en vivo & \cod{demo/alternativas/03-enjambre-juanpedro/GUION_3_MIN.md} \\
El pitch oficial y las preguntas del jurado & \cod{pitch/GUION.md}, \cod{pitch/PREGUNTAS.md} \\
\bottomrule
\end{tabularx}
\end{table}

\subsection{Comandos para comprobar lo principal}

Desde la raíz del repositorio, con Python 3.11 o superior:

\begin{consola}
pip install -r requirements.txt

# La suite completa
python -m pytest -q

# La demo, desde cero: 200 sin protección y 401 con Alicanto
python demo/run_split.py --out-dir demo/out

# El sello del set de prueba, recalculado sobre los archivos de hoy
python demo/alternativas/03-enjambre-juanpedro/hacer_evaluacion.py --sello

# La medición: 120 de 180 (unos 3 minutos)
python -m eval.final.proxy_flags --out /tmp/pf

# El pitch, con la parte en vivo
python demo/alternativas/03-enjambre-juanpedro/servidor.py
\end{consola}
